\documentclass[11pt, a4paper]{article}

% --- Packages ---
\usepackage[utf8]{inputenc}
\usepackage[T1]{fontenc}
\usepackage{mathptmx} % Times New Roman font for professional document typography
\usepackage{graphicx}
\usepackage[margin=0.8in]{geometry}
\usepackage{xcolor}
\usepackage{titlesec}
\usepackage{enumitem}
\usepackage{float}
\usepackage{booktabs}
\usepackage{multirow}
\usepackage{array}

% --- Amrita Colors ---
\definecolor{amritaMaroon}{HTML}{A4123F}

% --- Title & Section Formatting ---
\titlespacing*{\section}{0pt}{12pt}{4pt}
\titlespacing*{\subsection}{0pt}{8pt}{3pt}
\titleformat{\section}{\large\bfseries\color{amritaMaroon}}{\thesection}{0.8em}{}
\titleformat{\subsection}{\normalsize\bfseries\color{amritaMaroon}}{\thesubsection}{0.6em}{}

\pagestyle{plain}

\begin{document}

% ==========================================
% HEADER & TITLE BLOCK
% ==========================================
\noindent
\begin{minipage}{\textwidth}
    \centering
    {\Large \textbf{\textcolor{amritaMaroon}{
    AMRITA VISHWA VIDYAPEETHAM \\
    SCHOOL OF COMPUTING \\
    DEPARTMENT OF COMPUTER SCIENCE AND ENGINEERING \\
    }}}
    \vspace{0.2cm}
    {\large \textbf{[PROJECT TITLE HERE]}} \\[0.1cm]
    \textbf{23CSE498 -- Project Phase II (External Review)} \quad | \quad \textbf{Team Number: [GXXX]} \\
    \today
\end{minipage}

\vspace{0.3cm}
\hrule
\vspace{0.3cm}

% --- Team & Guide Details ---
\noindent
\begin{center}
    \renewcommand{\arraystretch}{1.45}
    \setlength{\tabcolsep}{6pt}
    \begin{tabular}{|l|p{4.8cm}|p{4.6cm}|p{3.2cm}|}
        \hline
        \multicolumn{2}{|c|}{\textbf{Team Members}} & \textbf{Project Guide} & \textbf{Project Type} \\
        \hline
        \small [CB.SC.U4CSE23XXX] & \small \raggedright [Student Name 1] &
        \multirow{4}{=}{\raggedright \small \textbf{[Dr/Mr/Ms. Guide Name]} \newline \textit{[Designation], \newline Dept. of CSE}} &
        \multirow{4}{=}{\raggedright \small \textbf{Type 2} \newline Application-Based Project} \\ \cline{1-2}
        \small [CB.SC.U4CSE23XXX] & \small [Student Name 2] & & \\ \cline{1-2}
        \small [CB.SC.U4CSE23XXX] & \small [Student Name 3] & & \\ \cline{1-2}
        \small [CB.SC.U4CSE23XXX] & \small [Student Name 4] & & \\
        \hline
    \end{tabular}
\end{center}
\vspace{0.2cm}

\section{Abstract}
\noindent
[Provide a concise 150--250 word summary of the entire project: state the real-world problem and target users, summarize the functional requirements and end-to-end application architecture, describe the complete (100\%) integrated implementation, and highlight the key outcomes from functional testing and prototype verification.]

% --- 2. PROBLEM DEFINITION & REQUIREMENT ANALYSIS ---
\section{Problem Definition \& Requirement Analysis}

\subsection{Problem Statement, Context \& Target Users}
\begin{itemize}[noitemsep, topsep=2pt]
    \item \textbf{Problem Definition:} {[Clearly define the real-world problem, operational bottleneck, or user need being addressed.]}
    \item \textbf{Target Users \& Stakeholders:} {[Identify the primary user groups and deployment context.]}
    \item \textbf{Existing Limitations:} {[Summarize limitations in existing tools or manual workflows \cite{ref1,ref2}.]}
    \item \textbf{Application Objectives:}
    \begin{enumerate}[noitemsep, topsep=2pt]
        \item {[Objective 1: State specific, measurable application objective]}
        \item {[Objective 2: State specific, measurable application objective]}
        \item {[Objective 3: State specific, measurable application objective]}
    \end{enumerate}
\end{itemize}

\subsection{Functional \& Non-Functional Requirements}
\begin{table}[H]
    \centering
    \footnotesize
    \renewcommand{\arraystretch}{1.25}
    \begin{tabular}{|c|p{3.0cm}|p{9.2cm}|}
        \hline
        \textbf{Req. ID} & \textbf{Category} & \textbf{Requirement Specification} \\
        \hline
        {[FR-01]} & Functional & {[Specify core user-facing feature or workflow requirement 1]} \\
        \hline
        {[FR-02]} & Functional & {[Specify core user-facing feature or workflow requirement 2]} \\
        \hline
        {[FR-03]} & Functional & {[Specify core user-facing feature or workflow requirement 3]} \\
        \hline
        {[NFR-01]} & Non-Functional & {[Specify latency, responsiveness, or throughput requirement]} \\
        \hline
        {[NFR-02]} & Non-Functional & {[Specify usability, reliability, safety, or hardware constraint]} \\
        \hline
    \end{tabular}
    \caption{Summary of Application Functional and Non-Functional Requirements}
\end{table}

% --- 3. APPLICATION ARCHITECTURE & SOLUTION DESIGN ---
\section{Application Architecture \& Solution Design}

\subsection{End-to-End Application Architecture}
\noindent
[Describe the overall system architecture and how user actions/inputs propagate across modules to generate outputs.]

\begin{figure}[H]
    \centering
    \fbox{\begin{minipage}[c][5.2cm]{0.95\textwidth}
        \centering
        \textbf{[Insert High-Resolution Application Architecture / End-to-End Workflow Diagram]} \\[0.2cm]
        \small {[Placeholder: Show User Interface, Core Processing Modules, Data Storage, and External/Hardware Interfaces]}
    \end{minipage}}
    \caption{End-to-End Application Architecture and Workflow}
\end{figure}

\subsection{Component Design, Data Flow \& Technology Stack}
\begin{itemize}[noitemsep, topsep=2pt]
    \item \textbf{Component Breakdown:} {[Describe each functional component and its responsibilities within the architecture.]}
    \item \textbf{User Interaction \& Interface Flow:} {[Explain the navigation and interaction design for end-users.]}
    \item \textbf{Data Storage \& External Interfaces:} {[Detail the data models/storage structure and any external APIs, sensors, or hardware interfaces.]}
    \item \textbf{Technology \& Deployment Stack:} {[Specify the selected platforms, frameworks, and runtime environment with technical justification.]}
\end{itemize}

% --- 4. COMPLETE IMPLEMENTATION DETAILS (100%) ---
\section{Implementation Details}

\subsection{Development \& Deployment Environment}
\begin{itemize}[noitemsep, topsep=2pt]
    \item \textbf{Platform \& Runtime Environment:} {[Specify operating environment, target devices, and deployment setup.]}
    \item \textbf{Libraries, Tools \& Hardware:} {[List software packages, SDKs, or hardware components used.]}
    \item \textbf{Source Code Repository:} {[Provide URL/reference to the complete, organized project repository.]}
\end{itemize}

\subsection{Module-Wise Implementation \& Integration Summary}
\begin{table}[H]
    \centering
    \footnotesize
    \renewcommand{\arraystretch}{1.25}
    \begin{tabular}{|p{3.5cm}|p{7.8cm}|c|}
        \hline
        \textbf{Application Module} & \textbf{Implemented Features \& Integration Details} & \textbf{Status} \\
        \hline
        {[Module 1: User Interface]} & {[Describe implemented screens, dashboards, forms, and user interactions]} & 100\% (Integrated) \\
        \hline
        {[Module 2: Core Logic]} & {[Describe implemented domain workflows, processing rules, and controllers]} & 100\% (Integrated) \\
        \hline
        {[Module 3: Data Layer]} & {[Describe implemented data storage, retrieval queries, and state persistence]} & 100\% (Integrated) \\
        \hline
        {[Module 4: External / IO]} & {[Describe integrated external APIs, hardware sensors/actuators, or services]} & 100\% (Functional) \\
        \hline
    \end{tabular}
    \caption{Application Module Implementation and Integration Status}
\end{table}

\subsection{Detailed Module Descriptions}
\begin{itemize}[noitemsep, topsep=2pt]
    \item \textbf{[Module 1 Name]:} {[Detail the implementation mechanics, inputs/outputs, and user workflow of Module 1.]}
    \item \textbf{[Module 2 Name]:} {[Detail the implementation mechanics, processing logic, and validation of Module 2.]}
    \item \textbf{[Module 3 Name]:} {[Detail the data handling, external communication, and integration of Module 3.]}
    \item \textbf{End-to-End Integration:} {[Explain how all modules are integrated into a unified, working application prototype.]}
\end{itemize}

% --- 5. FUNCTIONAL TESTING & PROTOTYPE DEMONSTRATION ---
\section{Functional Testing \& Prototype Demonstration}

\subsection{Functional Test Cases \& Verification}
\begin{table}[H]
    \centering
    \footnotesize
    \renewcommand{\arraystretch}{1.25}
    \begin{tabular}{|c|p{3.2cm}|p{3.4cm}|p{4.0cm}|c|}
        \hline
        \textbf{Test ID} & \textbf{Feature / Scenario} & \textbf{Input / User Action} & \textbf{Observed Output} & \textbf{Result} \\
        \hline
        {[TC-01]} & {[Primary User Flow]} & {[Specify test input]} & {[Matches expected behavior]} & {[Pass]} \\
        \hline
        {[TC-02]} & {[Secondary Feature]} & {[Specify test input]} & {[Matches expected behavior]} & {[Pass]} \\
        \hline
        {[TC-03]} & {[Module Integration]} & {[Specify test input]} & {[Matches expected behavior]} & {[Pass]} \\
        \hline
        {[TC-04]} & {[Boundary / Error Case]} & {[Invalid / edge input]} & {[Handled gracefully]} & {[Pass]} \\
        \hline
    \end{tabular}
    \caption{Functional Testing Summary Across Core Application Scenarios}
\end{table}

\subsection{Working Prototype Demonstration}
\begin{figure}[H]
    \centering
    \fbox{\begin{minipage}[c][4.8cm]{0.95\textwidth}
        \centering
        \textbf{[Insert Working Prototype Screenshots / Hardware Setup / End-to-End Execution Views]} \\[0.2cm]
        \small {[Placeholder: Include labeled screenshots demonstrating the working application from input to output]}
    \end{minipage}}
    \caption{Working Prototype Demonstration Snapshots}
\end{figure}

\noindent
\textbf{Evaluation \& Usability Summary:} {[Summarize operational performance, response time, accuracy, or user feedback observed during prototype verification.]}

% --- 6. TECHNICAL KNOWLEDGE & ARCHITECTURAL DECISIONS ---
\section{Technical Justification \& Architectural Decisions}
\begin{itemize}[noitemsep, topsep=2pt]
    \item \textbf{Technology Stack Rationale:} {[Justify why the chosen platforms, frameworks, and storage mechanisms are appropriate for this application.]}
    \item \textbf{Architectural \& Performance Trade-offs:} {[Explain trade-offs made regarding latency, resource consumption, modularity, and ease of deployment.]}
    \item \textbf{Security, Safety \& Error Handling:} {[Describe authentication, input validation, data privacy, and failure recovery mechanisms implemented.]}
\end{itemize}

% --- 7. CONCLUSION, LIMITATIONS & FUTURE SCOPE ---
\section{Conclusion, Limitations \& Future Scope}
\begin{itemize}[noitemsep, topsep=2pt]
    \item \textbf{Conclusion:} {[Summarize the completed application and how it fulfills the defined user requirements.]}
    \item \textbf{Current Operational Limitations:} {[State practical constraints of the current prototype.]}
    \item \textbf{Future Enhancements:} {[Outline planned features or scaling improvements.]}
\end{itemize}

% --- 8. DELIVERABLES & TECHNICAL ENRICHMENT SUMMARY ---
\section{Project Deliverables \& Technical Enrichment Summary}
\begin{table}[H]
    \centering
    \footnotesize
    \renewcommand{\arraystretch}{1.2}
    \begin{tabular}{|p{4.2cm}|p{7.5cm}|c|}
        \hline
        \textbf{Deliverable Category} & \textbf{Details / Title / Repository Link} & \textbf{Status} \\
        \hline
        Source Code Repository & {[Repository URL and documentation status]} & {[Completed]} \\
        \hline
        Individual Online Certifications & {[List completed 10--20 hr courses for each team member]} & {[Completed]} \\
        \hline
        Group TEA Deliverable 1 & {[Title and details of completed Group Activity 1]} & {[Completed]} \\
        \hline
        Group TEA Deliverable 2 & {[Title and details of completed Group Activity 2]} & {[Completed]} \\
        \hline
        Group TEA Deliverable 3 & {[Title and details of completed Group Activity 3]} & {[Completed]} \\
        \hline
    \end{tabular}
    \caption{Summary of Completed Project and Technical Enrichment Deliverables}
\end{table}

% --- REFERENCES ---
\vspace{0.2cm}
\nocite{*}
\renewcommand{\refname}{\large\bfseries\color{amritaMaroon}References}
\bibliographystyle{ieeetr}
\bibliography{ref}

\end{document}
